\documentclass[10pt,a4paper]{amsart}
\usepackage{amsmath,amssymb,amsthm,graphicx}
% Retained minimal preamble for the Stacks Project chapter flat.tex (More on Flatness).
% Verbatim excerpt of the chapter preamble: Xy-pic setup (native preamble lines 14-18) and the native
% theorem declarations with the mathematical macros (native preamble lines 154-217). Copied mechanically;
% nothing added, removed or redefined. The wrapper class and the amsmath/amssymb/amsthm/graphicx base
% packages are supplied by the assembler, so the class/externaldocument part of the native preamble is omitted.
% For commutative diagrams we use Xy-pic
\usepackage[all]{xy}

% We use 2cell for 2-commutative diagrams.
\xyoption{2cell}
% Theorem environments.
%
\theoremstyle{plain}
\newtheorem{theorem}[subsection]{Theorem}
\newtheorem{proposition}[subsection]{Proposition}
\newtheorem{lemma}[subsection]{Lemma}

\theoremstyle{definition}
\newtheorem{definition}[subsection]{Definition}
\newtheorem{example}[subsection]{Example}
\newtheorem{exercise}[subsection]{Exercise}
\newtheorem{situation}[subsection]{Situation}

\theoremstyle{remark}
\newtheorem{remark}[subsection]{Remark}
\newtheorem{remarks}[subsection]{Remarks}

\numberwithin{equation}{subsection}

% Macros
%
\def\lim{\mathop{\mathrm{lim}}\nolimits}
\def\colim{\mathop{\mathrm{colim}}\nolimits}
\def\Spec{\mathop{\mathrm{Spec}}}
\def\Hom{\mathop{\mathrm{Hom}}\nolimits}
\def\Ext{\mathop{\mathrm{Ext}}\nolimits}
\def\SheafHom{\mathop{\mathcal{H}\!\mathit{om}}\nolimits}
\def\SheafExt{\mathop{\mathcal{E}\!\mathit{xt}}\nolimits}
\def\Sch{\mathit{Sch}}
\def\Mor{\mathop{\mathrm{Mor}}\nolimits}
\def\Ob{\mathop{\mathrm{Ob}}\nolimits}
\def\Sh{\mathop{\mathit{Sh}}\nolimits}
\def\NL{\mathop{N\!L}\nolimits}
\def\CH{\mathop{\mathrm{CH}}\nolimits}
\def\proetale{{pro\text{-}\acute{e}tale}}
\def\etale{{\acute{e}tale}}
\def\QCoh{\mathit{QCoh}}
\def\Ker{\mathop{\mathrm{Ker}}}
\def\Im{\mathop{\mathrm{Im}}}
\def\Coker{\mathop{\mathrm{Coker}}}
\def\Coim{\mathop{\mathrm{Coim}}}

% Boxtimes
%
\DeclareMathSymbol{\boxtimes}{\mathbin}{AMSa}{"02}

%
% Macros for moduli stacks/spaces
%
\def\QCohstack{\mathcal{QC}\!\mathit{oh}}
\def\Cohstack{\mathcal{C}\!\mathit{oh}}
\def\Spacesstack{\mathcal{S}\!\mathit{paces}}
\def\Quotfunctor{\mathrm{Quot}}
\def\Hilbfunctor{\mathrm{Hilb}}
\def\Curvesstack{\mathcal{C}\!\mathit{urves}}
\def\Polarizedstack{\mathcal{P}\!\mathit{olarized}}
\def\Complexesstack{\mathcal{C}\!\mathit{omplexes}}
% \Pic is the operator that assigns to X its picard group, usage \Pic(X)
% \Picardstack_{X/B} denotes the Picard stack of X over B
% \Picardfunctor_{X/B} denotes the Picard functor of X over B
\def\Pic{\mathop{\mathrm{Pic}}\nolimits}
\def\Picardstack{\mathcal{P}\!\mathit{ic}}
\def\Picardfunctor{\mathrm{Pic}}
\def\Deformationcategory{\mathcal{D}\!\mathit{ef}}

% --- appended 2026-10-12 (Chief): the project own definitions of the reference/slogan/history
% display environments, copied mechanically from the shared Stacks preamble lines 8-12
% (\usepackage{verbatim} plus the three \newenvironment lines, which the chapter preamble keeps
% verbatim from the project preamble). Nothing added, removed or redefined: these environments
% carry no mathematical content and the project itself typesets them as comments. Without these
% three lines a counted statement whose own text uses them fails with Environment slogan undefined.
\usepackage{verbatim}
\newenvironment{reference}{\comment}{\endcomment}
\newenvironment{slogan}{\comment}{\endcomment}
\newenvironment{history}{\comment}{\endcomment}
\usepackage{cleveref}
\crefname{originallemma}{original lemma}{original lemmas}
\Crefname{originallemma}{Original lemma}{Original lemmas}
\makeatletter
\newlabel{algebra-lemma-NAK}{{00DV}{source}{}{sourcepointer-1}{}}
\newlabel{algebra-lemma-NAK@cref}{{[originallemma][0][]00DV}{[1][0][]source}{}{}{}}
\makeatother
\crefname{originallemma}{original lemma}{original lemmas}
\Crefname{originallemma}{Original lemma}{Original lemmas}
\makeatletter
\newlabel{algebra-lemma-flat-ses}{{00HM}{source}{}{sourcepointer-2}{}}
\newlabel{algebra-lemma-flat-ses@cref}{{[originallemma][0][]00HM}{[1][0][]source}{}{}{}}
\makeatother
\crefname{originallemma}{original lemma}{original lemmas}
\Crefname{originallemma}{Original lemma}{Original lemmas}
\makeatletter
\newlabel{algebra-lemma-ui-flat-domain}{{058P}{source}{}{sourcepointer-3}{}}
\newlabel{algebra-lemma-ui-flat-domain@cref}{{[originallemma][0][]058P}{[1][0][]source}{}{}{}}
\makeatother
\crefname{originallemma}{original lemma}{original lemmas}
\Crefname{originallemma}{Original lemma}{Original lemmas}
\makeatletter
\newlabel{algebra-lemma-universally-injective-localize}{{05CM}{source}{}{sourcepointer-4}{}}
\newlabel{algebra-lemma-universally-injective-localize@cref}{{[originallemma][0][]05CM}{[1][0][]source}{}{}{}}
\makeatother
\crefname{originallemma}{original lemma}{original lemmas}
\Crefname{originallemma}{Original lemma}{Original lemmas}
\makeatletter
\newlabel{algebra-lemma-universally-injective-permanence}{{05CJ}{source}{}{sourcepointer-5}{}}
\newlabel{algebra-lemma-universally-injective-permanence@cref}{{[originallemma][0][]05CJ}{[1][0][]source}{}{}{}}
\makeatother
\crefname{originallemma}{original lemma}{original lemmas}
\Crefname{originallemma}{Original lemma}{Original lemmas}
\makeatletter
\newlabel{lemma-fibres-irreducible-flat-projective-nonnoetherian}{{05FT}{source}{}{sourcepointer-6}{}}
\newlabel{lemma-fibres-irreducible-flat-projective-nonnoetherian@cref}{{[originallemma][0][]05FT}{[1][0][]source}{}{}{}}
\makeatother
\crefname{originallemma}{original lemma}{original lemmas}
\Crefname{originallemma}{Original lemma}{Original lemmas}
\makeatletter
\newlabel{lemma-invert-universally-injective}{{05FR}{source}{}{sourcepointer-7}{}}
\newlabel{lemma-invert-universally-injective@cref}{{[originallemma][0][]05FR}{[1][0][]source}{}{}{}}
\makeatother
\crefname{originallemma}{original lemma}{original lemmas}
\Crefname{originallemma}{Original lemma}{Original lemmas}
\makeatletter
\newlabel{lemma-universally-injective-local}{{05FQ}{source}{}{sourcepointer-8}{}}
\newlabel{lemma-universally-injective-local@cref}{{[originallemma][0][]05FQ}{[1][0][]source}{}{}{}}
\makeatother
\title{The algebraic induction step: flatness of a localised module over a flat finitely presented ring map with geometrically integral fibres}
\author{The Stacks Project Authors}
\date{stacks-project:flat.tex:local-snapshot-20261011; GNU FDL 1.2 or later}
\begin{document}
\maketitle
\noindent\textit{Editorial note.} All delivered mathematical text is the author's own, copied from the pinned Stacks Project chapter file flat.tex; no statement or proof step was rewritten, reordered or removed.\par\medskip
\noindent\textit{Editorial note.} Preamble provenance: the chapter's own retained minimal preamble (sources/stacks-flat/retained-preamble.tex, lines 1-85) is a verbatim excerpt of the native Xy-pic setup and of the native theorem declarations and mathematical macros, followed by the project's own reference/slogan/history display-environment definitions copied mechanically from the shared preamble; the wrapper class and the amsmath/amssymb/amsthm/graphicx base packages are supplied by the assembler, and extra\_packages is left empty because everything needed is already in the retained preamble. No macro was invented or redefined.\par\medskip
\noindent\textit{Editorial note.} Same-chapter pointers are located at their \textbackslash{}label lines inside the pinned chapter file; cross-chapter pointers are located in the project's own tag list staged in the snapshot (sources/stacks-chow/tags.txt), at the verbatim lines carrying their canonical tags.\par\medskip
\noindent\textit{Allowed standard prerequisites (editorial).}\begin{itemize}
\item Universal injectivity over a local base for a map whose reduction modulo the maximal ideal is injective. Same chapter flat.tex, labels lemma-universally-injective-local (proof line 2396) and lemma-fibres-irreducible-flat-projective-nonnoetherian (proof line 2397). Together they show that the map from the free module to the localised module is universally injective over R once its reduction modulo the maximal ideal is injective, which is the heart of the implication (1) implies (2).
\item Permanence and localisation of universal injectivity, and flatness of the cokernel of a universally injective map over a domain. Algebra, labels algebra-lemma-universally-injective-permanence (proof line 2400), algebra-lemma-universally-injective-localize (proof line 2403) and algebra-lemma-ui-flat-domain (proof line 2405). Carries universal injectivity from the localised map back to the map on the local ring and makes the cokernel flat over R, which produces condition (2).
\item Nakayama's lemma, nonzerodivisors of a flat finitely presented ring map, and flatness of an extension. Algebra, labels algebra-lemma-NAK (proof line 2411); Same chapter flat.tex, labels lemma-invert-universally-injective (proof line 2417) and lemma-fibres-irreducible-flat-projective-nonnoetherian (proof line 2418); Algebra, label algebra-lemma-flat-ses (proof line 2421). Nakayama upgrades injectivity at p to an isomorphism, the nonzerodivisor statement lets injectivity at p be pulled back to injectivity at q, and flatness of an extension closes the cycle from (5) back to (1).
\end{itemize}\medskip
\ifdefined\hypertarget\hypertarget{sourcepointer-1}{}\fi
\noindent\textit{Source reference.} Original lemma 00DV: sources/stacks-chow/tags.txt, source lines 512--512. Navigation pointer for \textbackslash{}ref\{algebra-lemma-NAK\} used inside the counted span (line 2411); the argument that consumes it is retained in the span.\par\medskip
\ifdefined\hypertarget\hypertarget{sourcepointer-2}{}\fi
\noindent\textit{Source reference.} Original lemma 00HM: sources/stacks-chow/tags.txt, source lines 648--648. Navigation pointer for \textbackslash{}ref\{algebra-lemma-flat-ses\} used inside the counted span (line 2421); the argument that consumes it is retained in the span.\par\medskip
\ifdefined\hypertarget\hypertarget{sourcepointer-3}{}\fi
\noindent\textit{Source reference.} Original lemma 058P: sources/stacks-chow/tags.txt, source lines 6787--6787. Navigation pointer for \textbackslash{}ref\{algebra-lemma-ui-flat-domain\} used inside the counted span (line 2405); the argument that consumes it is retained in the span.\par\medskip
\ifdefined\hypertarget\hypertarget{sourcepointer-4}{}\fi
\noindent\textit{Source reference.} Original lemma 05CM: sources/stacks-chow/tags.txt, source lines 6925--6925. Navigation pointer for \textbackslash{}ref\{algebra-lemma-universally-injective-localize\} used inside the counted span (line 2403); the argument that consumes it is retained in the span.\par\medskip
\ifdefined\hypertarget\hypertarget{sourcepointer-5}{}\fi
\noindent\textit{Source reference.} Original lemma 05CJ: sources/stacks-chow/tags.txt, source lines 6922--6922. Navigation pointer for \textbackslash{}ref\{algebra-lemma-universally-injective-permanence\} used inside the counted span (line 2400); the argument that consumes it is retained in the span.\par\medskip
\ifdefined\hypertarget\hypertarget{sourcepointer-6}{}\fi
\noindent\textit{Source reference.} Original lemma 05FT: sources/stacks-flat/pinned-original-source.tex, source lines 2247--2247. Navigation pointer for \textbackslash{}ref\{lemma-fibres-irreducible-flat-projective-nonnoetherian\} used inside the counted span (line 2397); the argument that consumes it is retained in the span.\par\medskip
\ifdefined\hypertarget\hypertarget{sourcepointer-7}{}\fi
\noindent\textit{Source reference.} Original lemma 05FR: sources/stacks-flat/pinned-original-source.tex, source lines 1976--1976. Navigation pointer for \textbackslash{}ref\{lemma-invert-universally-injective\} used inside the counted span (line 2417); the argument that consumes it is retained in the span.\par\medskip
\ifdefined\hypertarget\hypertarget{sourcepointer-8}{}\fi
\noindent\textit{Source reference.} Original lemma 05FQ: sources/stacks-flat/pinned-original-source.tex, source lines 1946--1946. Navigation pointer for \textbackslash{}ref\{lemma-universally-injective-local\} used inside the counted span (line 2396); the argument that consumes it is retained in the span.\par\medskip
\begin{lemma}
\label{lemma-induction-step}
Let $(R, \mathfrak m)$ be a local ring. Let $R \to S$ be a finitely presented
flat ring map with geometrically integral fibres. Write
$\mathfrak p = \mathfrak mS$. Let $\mathfrak q \subset S$ be a prime ideal
lying over $\mathfrak m$. Let $N$ be a finite $S$-module.
There exist $r \geq 0$ and an $S$-module map
$$
\alpha : S^{\oplus r} \longrightarrow N
$$
such that
$\alpha : \kappa(\mathfrak p)^{\oplus r} \to N \otimes_S \kappa(\mathfrak p)$
is an isomorphism. For any such $\alpha$ the following are equivalent:
\begin{enumerate}
\item $N_{\mathfrak q}$ is $R$-flat,
\item $\alpha$ is $R$-universally injective and
$\Coker(\alpha)_{\mathfrak q}$ is $R$-flat,
\item $\alpha$ is injective and
$\Coker(\alpha)_{\mathfrak q}$ is $R$-flat,
\item $\alpha_{\mathfrak p}$ is an isomorphism and
$\Coker(\alpha)_{\mathfrak q}$ is $R$-flat, and
\item $\alpha_{\mathfrak q}$ is injective and
$\Coker(\alpha)_{\mathfrak q}$ is $R$-flat.
\end{enumerate}
\end{lemma}

\begin{proof}
To obtain $\alpha$ set
$r = \dim_{\kappa(\mathfrak p)} N \otimes_S \kappa(\mathfrak p)$ and pick
$x_1, \ldots, x_r \in N$ which form a basis of
$N \otimes_S \kappa(\mathfrak p)$. Define
$\alpha(s_1, \ldots, s_r) = \sum s_i x_i$. This proves the existence.

\medskip\noindent
Fix an $\alpha$. The most interesting implication is
(1) $\Rightarrow$ (2) which we prove first. Assume (1).
Because $S/\mathfrak mS$ is a domain with fraction field $\kappa(\mathfrak p)$
we see that
$(S/\mathfrak mS)^{\oplus r} \to
N_{\mathfrak p}/\mathfrak mN_{\mathfrak p} = N \otimes_S \kappa(\mathfrak p)$
is injective. Hence by
Lemmas \ref{lemma-universally-injective-local} and
\ref{lemma-fibres-irreducible-flat-projective-nonnoetherian}.
the map $S^{\oplus r} \to N_{\mathfrak p}$ is $R$-universally injective.
It follows that $S^{\oplus r} \to N$ is $R$-universally injective, see
Algebra, Lemma \ref{algebra-lemma-universally-injective-permanence}.
Then also the localization $\alpha_{\mathfrak q}$ is $R$-universally
injective, see
Algebra, Lemma \ref{algebra-lemma-universally-injective-localize}.
We conclude that $\Coker(\alpha)_{\mathfrak q}$ is $R$-flat by
Algebra, Lemma \ref{algebra-lemma-ui-flat-domain}.

\medskip\noindent
The implication (2) $\Rightarrow$ (3) is immediate. If (3) holds, then
$\alpha_{\mathfrak p}$ is injective as a localization of an injective
module map. By Nakayama's lemma
(Algebra, Lemma \ref{algebra-lemma-NAK})
$\alpha_{\mathfrak p}$ is surjective too. Hence (3) $\Rightarrow$ (4).
If (4) holds, then $\alpha_{\mathfrak p}$ is an isomorphism, so
$\alpha$ is injective as $S_{\mathfrak q} \to S_{\mathfrak p}$ is injective.
Namely, elements of $S \setminus \mathfrak p$ are nonzerodivisors on $S$
by a combination of
Lemmas \ref{lemma-invert-universally-injective} and
\ref{lemma-fibres-irreducible-flat-projective-nonnoetherian}.
Hence (4) $\Rightarrow$ (5). Finally, if (5) holds, then
$N_{\mathfrak q}$ is $R$-flat as an extension of flat modules, see
Algebra, Lemma \ref{algebra-lemma-flat-ses}.
Hence (5) $\Rightarrow$ (1) and the proof is finished.
\end{proof}

\end{document}
